\section{Introduction}\label{sec:intro}

Teaching is generally regarded as an emotionally demanding occupation, because
teachers work with many people at once and are responsible for their results.
Such conditions are associated with professional burnout, a state that Maslach
and colleagues describe through three dimensions: emotional exhaustion,
depersonalization and reduced personal accomplishment~\cite{maslach1981mbi,maslach2001job}.
A systematic review and meta-analysis of 45 articles on secondary school
teachers concluded that this group is at high risk of burnout~\cite{garciacarmona2019burnout}.

The risk is not necessarily the same in every school. The job
demands--resources model explains burnout by the balance between job demands
and the resources available to cope with them~\cite{demerouti2001jdr,bakker2007jdr};
it is flexible and can be applied to a wide range of occupations~\cite{bakker2007jdr}.
State and private schools may differ in exactly these conditions, for example
in pay, class size, administrative control and autonomy.

Comparisons of the two school types exist, but they do not give a clear
picture. In a study of special-school teachers (120 per school type), burnout was
higher in government schools than in private schools~\cite{doss2018special},
whereas in a study of 180 teachers with the MBI, teachers of English in Iranian
private language institutes reported worse burnout than their colleagues in
public schools~\cite{bahrami2019efl}. These
studies concern specific groups, they point in different directions, and the
review cited above shows that the measured prevalence depends on the
instrument~\cite{garciacarmona2019burnout}, which makes studies hard to
compare. We therefore plan a direct comparison of teachers in state and
private schools with one instrument and one analysis protocol.
% TODO (check before submission): search once more for studies of general-education teachers in our region.

We describe a planned cross-sectional study. Burnout is measured with the
Maslach Burnout Inventory in its Russian-language adaptation~\cite{vodopyanova2008burnout},
and the groups are compared with Welch's $t$-test, the Mann--Whitney test and
bootstrap confidence intervals. The study addresses three research questions:
RQ1, do burnout levels of the two groups differ (main comparison: emotional
exhaustion); RQ2, do the groups differ on each of the three scales; and RQ3,
which working conditions are associated with burnout (exploratory). The
analysis runs in a containerised pipeline, which supports reproducibility~\cite{boettiger2015docker}.

Questionnaire data may arrive in several waves, and every wave requires the
scores and tests to be recomputed. If the analysis is slow or cannot be repeated
exactly, errors and undocumented manual steps creep in. We therefore treat the
pipeline as a part of the study and check how its latency and throughput change
when the number of respondents grows.

The contributions of this work are:
\begin{itemize}
  \item a research design and analysis protocol for comparing teacher burnout
        in state and private schools, fixed before data collection;
  \item three research questions with testable hypotheses (RQ1--RQ3);
  \item an open, reproducible software pipeline for scoring the questionnaire
        and computing the statistics;
  \item a measurement of the pipeline's latency and throughput that reports
        tail percentiles~\cite{dean2013tail} with bootstrap confidence intervals.
\end{itemize}

The rest of the paper is organised as follows. \autoref{sec:method} describes
the design (\autoref{sec:design}), the instrument (\autoref{sec:instrument}), the
statistical methods (\autoref{sec:tests}), the pipeline (\autoref{sec:pipeline}) and
the measurement procedure (\autoref{sec:measure}). \autoref{sec:results} gives the
setup (\autoref{sec:setup}), the results (\autoref{sec:res}) and the discussion
(\autoref{sec:disc}), and \autoref{sec:concl} states the limitations
(\autoref{sec:limits}), concludes and outlines future work (\autoref{sec:future}).
